% Source: https://artofproblemsolving.com/community/c3h3040907_circle_chord
% Stable identity: aops-topic-3040907-q01
If the inner radius is $r$, then $\pi(100-r^2)=64\pi$, giving $r=6$. The perpendicular distance from the common centre to the chord is 6. The radius 10 and half-chord form a right triangle, so the half-chord is $\sqrt{100-36}=8$. The chord length is $\boxed{16}$.
\begin{asy}
usepackage("amsmath");
size(7cm);
draw(circle((0,0),10));draw(circle((0,0),6));draw((-8,6)--(8,6),blue);draw((0,0)--(0,6)--(8,6)--cycle,dashed);label("$6$",(0,3),W);label("$10$",(4,3),SE);
\end{asy}
